\section{Methodology}
\label{sec:method}

\subsection{The universality boundary and granularity mismatch}

A decode-time intervention module fails to be strictly universal if either its inputs or outputs are host-specific. Earlier calibration methods break this universality boundary by consuming model-native hidden states and projecting them directly into the host model's lexical geometry. Such learned mappings are host-coupled by construction and cannot be transferred across heterogeneous Multimodal Large Language Models (MLLMs) without retraining.

To achieve true hot-swappability, an intervention mechanism must operate in a standardized external semantic space. However, designing such an external controller requires addressing the fundamental root cause of hallucinations: the \textbf{cross-modal granularity mismatch}. While MLLMs encode images as continuous, holistic feature maps, autoregressive generation forces them to output discrete, often semantically fragmented sub-word tokens (e.g., \texttt{["gi", "\#\#raf", "\#\#fe"]}). It is mathematically intractable to establish reliable physical grounding between a continuous visual entity and a meaningless, fragmented token. This granularity collapse forces the MLLM into an \emph{open-loop decoding state} where visual verification fails, allowing the statistical inertia of language priors to overwhelm visual constraints. CHORD is designed to bridge this mismatch and forcefully close the decoding loop.

\subsection{The CHORD framework overview}

CHORD (Calibrating Hallucinations via Object-Resonant Decoding) is a training-free, plug-and-play decoding framework. Instead of overriding the full vocabulary distribution, CHORD acts exclusively on the host model's active decoding frontier
\begin{equation}
  \mathcal{F}_t = \mathrm{Top}\text{-}K(L_t),
\end{equation}
where $L_t$ represents the original logits at step $t$. The framework operates through four sequential stages: (1)~Tokenizer Bridging, (2)~Continuous Resonance Monitoring, (3)~Hypothesis-Driven Resonance Delta, and (4)~Logit Calibration.

Before generation begins, CHORD performs a one-time offline extraction. Given the input image $I$, a frozen external vision encoder extracts a dense local region feature bank $\mathcal{Z}_R = \{z_{R_1}, z_{R_2}, \dots, z_{R_N}\}$. This bank serves as a static, objective physical reality cache for subsequent resonance verification.

\subsection{Parameter-free tokenizer bridge}

To overcome the granularity mismatch, CHORD does not evaluate raw token IDs. Let $B_m$ denote a parameter-free tokenizer bridge associated with the host model. For each candidate token $c_k \in \mathcal{F}_t$, the bridge performs a deterministic lookahead (utilizing the tokenizer's prefix tree) to lift the token into a semantically complete text span $A_k$.

This yields a set of candidate spans $\mathcal{S}_t = \{A_1, A_2, \dots, A_K\}$. This bridge introduces no learned parameters. It is the critical mechanism that transforms model-specific, fragmented token proposals into cross-model, physically groundable semantic propositions.

\subsection{Hypothesis-driven resonance delta}

Instead of training a black-box MLP to blindly score text-image pairs, CHORD models decoding as a physical resonance process. We treat the text prefix $p$ appended with a candidate span $A_k$ as a ``semantic sonar probe'' directed at the visual region bank $\mathcal{Z}_R$.

To strictly prevent \emph{spatial misalignment} (where the prefix and the candidate match different, unrelated regions), CHORD calculates the resonance shift anchored to the exact same physical region. First, we identify the single most relevant region $r^*$ for the combined hypothesis:
\begin{equation}
  r^* = \arg\max_{r \in \mathcal{Z}_R} \cos\bigl(E_{\text{text}}(p \oplus A_k), z_r\bigr),
\end{equation}
where $E_{\text{text}}$ is a frozen public text encoder. Next, we compute the \textbf{Resonance Delta ($\Delta S$)}, which isolates the physical impact of adding candidate $A_k$:
\begin{equation}
  \Delta S(A_k) = \cos\bigl(E_{\text{text}}(p \oplus A_k), z_{r^*}\bigr) - \cos\bigl(E_{\text{text}}(p), z_{r^*}\bigr).
\end{equation}

This formulation acts as a \textbf{heuristic-free filter}:
\begin{enumerate}
  \item \textbf{True visual spans ($\Delta S > 0$):} If $A_k$ is a visually supported object or correct spatial relation (e.g., ``left''), it increases the semantic alignment with region $r^*$, yielding a positive gradient.
  \item \textbf{Hallucinations ($\Delta S < 0$):} If $A_k$ contradicts the image (e.g., inventing an object), the alignment collapses, yielding a negative gradient.
  \item \textbf{Non-visual words ($\Delta S \approx 0$):} If $A_k$ is an abstract or grammatical word (e.g., ``very'', ``is''), $\Delta S \approx 0$. Thus, CHORD safely ignores non-visual words without relying on brittle POS taggers.
\end{enumerate}

\subsection{Overcoming overconfidence: continuous monitoring and system optimization}

A naive entropy-based gate fails to detect \emph{overconfidence hallucinations}, where strong language priors cause the model to output hallucinations with extremely low entropy. To address this while maintaining near-zero inference overhead, CHORD abandons entropy gating in favor of \textbf{Continuous Top-$1$ Monitoring with Full-$K$ Backoff}, powered by Key-Value (KV) cache acceleration.

Because the external $E_{\text{text}}$ utilizes causal masked self-attention, the prefix $p$ is continuously KV-cached. Extracting the feature for $p \oplus A_k$ requires only 1--2 tokens of forward propagation. CHORD continuously calculates $\Delta S(A_{\text{top1}})$ for the host model's most confident token at every step:
\begin{itemize}
  \item If $\Delta S(A_{\text{top1}}) \ge 0$, the token is visually safe or irrelevant, and CHORD allows generation to proceed.
  \item If $\Delta S(A_{\text{top1}}) < 0$, a language-prior-driven hallucination is detected regardless of confidence. Only then does CHORD trigger full evaluation over all $K$ candidates to find a grounded alternative.
\end{itemize}
This system-level optimization reduces computational overhead while rigorously preventing overconfidence failures.

\subsection{Logit calibration}

When full intervention is triggered, CHORD calibrates the active frontier by injecting resonance evidence back into the host model's logits. The adjusted logit for candidate token $c_k$ (corresponding to span $A_k$) is
\begin{equation}
  \tilde{L}_t(c_k) = L_t(c_k) + \alpha \cdot \Delta S(A_k),
\end{equation}
where $\alpha$ is a scaling hyperparameter. The logits are then re-normalized via softmax. Through this bounded redistribution, CHORD nudges the host model's trajectory toward visual reality without destroying native linguistic structures or requiring parameter updates.
